\documentclass[tikz,border=4pt]{standalone}
% ==========================================================================
% tikz-tcee.tex — Estilos comunes de los gráficos TikZ del temario
% Autor: Víctor Gutiérrez Marcos
%
% Cada gráfico es un fichero standalone en fuentes/<tema>/graficos/:
%
%   \documentclass[tikz,border=4pt]{standalone}
%   % ==========================================================================
% tikz-tcee.tex — Estilos comunes de los gráficos TikZ del temario
% Autor: Víctor Gutiérrez Marcos
%
% Cada gráfico es un fichero standalone en fuentes/<tema>/graficos/:
%
%   \documentclass[tikz,border=4pt]{standalone}
%   % ==========================================================================
% tikz-tcee.tex — Estilos comunes de los gráficos TikZ del temario
% Autor: Víctor Gutiérrez Marcos
%
% Cada gráfico es un fichero standalone en fuentes/<tema>/graficos/:
%
%   \documentclass[tikz,border=4pt]{standalone}
%   \input{../../../latex/tikz-tcee}
%   \begin{document}
%   \begin{tikzpicture}[tcee]
%     \ejes{Q}{P}{6}{5}
%     \draw[curva1] (0.5,4.5) -- (5,0.8) node[etiqueta,right] {$D$};
%     \capa{2}{\draw[curva2] ...;}   % en el vídeo aparece en el paso 2
%   \end{tikzpicture}
%   \end{document}
%
% Paleta: la de la web (morado, dorado, salvia) más tres colores de apoyo
% distinguibles en blanco y negro por el trazo.
% ==========================================================================
\usepackage[T1]{fontenc}
\usepackage{newpxtext,newpxmath}
\usepackage{xcolor}
\usepackage{tikz,pgfplots}
\pgfplotsset{compat=1.18}
\usetikzlibrary{arrows.meta,calc,positioning,patterns,decorations.pathreplacing,intersections,shapes.misc,backgrounds}

\definecolor{tceemorado}{HTML}{5F2987}
\definecolor{tceedorado}{HTML}{B8860B}
\definecolor{tceeverde}{HTML}{3C7D3A}
\definecolor{tceeazul}{HTML}{1F5F99}
\definecolor{tceerojo}{HTML}{B22222}
\definecolor{tceenaranja}{HTML}{E07B28}
\definecolor{tceegris}{HTML}{767171}
\definecolor{tceesalvia}{HTML}{E2EFD9}

% Capas para el vídeo: \capa{n}{…} dibuja su contenido solo a partir del paso n.
% En el PDF, la web y Word se ve todo; generar-video.py compila el gráfico una
% vez por paso con \def\tceepaso{k} para que cada elemento aparezca al narrarlo.
% Vectores con flecha encima, como en el texto (tcee.sty): \vect{x}, \vect{p}
\providecommand{\vect}[1]{\vec{#1}}
\providecommand{\tceepaso}{1000}
\newcommand{\capa}[2]{\ifnum#1>\tceepaso\relax\else#2\fi}
% Gráficos por pasos del mismo tamaño (propuesta del 3.A.29, 11/11/2026): con
% \def\tceefantasma{1} antes del \input (p. ej., x-paso1.tex:
% «\def\tceepaso{3}\def\tceefantasma{1}\input{x.tex}»), las capas posteriores
% al paso se dibujan invisibles y ocupan su sitio. Sin \tceefantasma (el vídeo),
% no cambia nada.
\ifdefined\tceefantasma
  \renewcommand{\capa}[2]{\ifnum#1>\tceepaso\relax\begin{scope}[opacity=0]#2\end{scope}\else#2\fi}
\fi

\tikzset{
  tcee/.style={line cap=round,line join=round,font=\small,>={Stealth[length=2.2mm]}},
  eje/.style={->,thick,black},
  curva1/.style={very thick,tceemorado},
  curva2/.style={very thick,tceeazul},
  curva3/.style={very thick,tceeverde},
  curva4/.style={very thick,tceedorado},
  curvanueva/.style={very thick,tceerojo},
  desplazada/.style={very thick,dashed},
  guia/.style={thin,dashed,tceegris},
  punto/.style={circle,fill=black,inner sep=1.3pt},
  etiqueta/.style={font=\small,inner sep=2pt},
  flecha/.style={->,thick,tceerojo},
  area/.style={fill=tceemorado,fill opacity=0.18,draw=none},
  area2/.style={fill=tceedorado,fill opacity=0.22,draw=none},
}

% \ejes{etiqueta x}{etiqueta y}{largo x}{alto y}
\newcommand{\ejes}[4]{%
  \draw[eje] (0,0) -- (#3,0) node[below left=2pt and 0pt] {#1};
  \draw[eje] (0,0) -- (0,#4) node[left] {#2};
  \node[below left] at (0,0) {$0$};}
% \proyeccion{(x,y)}{etq x}{etq y}: líneas guía a los ejes con etiquetas
\newcommand{\proyeccion}[3]{%
  \draw[guia] let \p1=#1 in (\x1,0) node[below] {#2} |- (0,\y1) node[left] {#3};}

% --------------------------------------------------------------------------
% Paleta «dictamen» (dic…): gráficos de datos de la Guía del dictamen (1D01) y
% de cualquier gráfico de coyuntura que quiera los mismos colores. Un concepto,
% un color en todos los gráficos (CRITERIOS §7, «tres reglas», 24/10/2026);
% tabla completa en fuentes/1D01/_trabajo/paleta-1D01.md. No cambia nada de
% lo anterior: son colores y estilos nuevos.
% Actividad: el PIB (y sus variantes: PIB por persona, PIB mundial), morado.
\definecolor{dicpib}{HTML}{5F2987}
% Demanda interna, gama cálida: consumo (naranjas: privado oscuro, público
% claro), inversión (verde) y existencias (gris con trama, aparte).
\definecolor{dicconspriv}{HTML}{C76A1F}
\definecolor{dicconspub}{HTML}{F5C690}
\definecolor{dicfbcf}{HTML}{2A6B3C}
\definecolor{dicexistencias}{HTML}{B4B4B4}
% Demanda externa y comercio, gama fría (azul claro, que se separa del morado)
\definecolor{dicexterna}{HTML}{6AA6DA}
% Zonas: zona del euro (azul marino), España (rojo), resto de países (grises)
\definecolor{dicze}{HTML}{17406D}
\definecolor{dicespana}{HTML}{C8382E}
\definecolor{dicotro}{HTML}{8A8A8A}
\definecolor{dicotroclaro}{HTML}{B5B5B5}
\definecolor{dicotromasclaro}{HTML}{DADADA}
% Precios: inflación general (negro, continuo) y subyacente (negro,
% discontinuo); componentes: energía (bermellón), alimentos (amarillo) y los
% de la subyacente en azules (bienes industriales claro, servicios oscuro).
\definecolor{dicinflacion}{HTML}{1A1A1A}
\definecolor{dicenergia}{HTML}{D55E00}
\definecolor{dicalimentos}{HTML}{F2D45C}
\definecolor{dicbine}{HTML}{7FB2E0}
\definecolor{dicservicios}{HTML}{1F4E89}
% Trabajo: empleo (verde azulado), horas (el mismo, claro), tasa de actividad
% (verde azulado oscuro), paro (teja); productividad (dorado de la web).
\definecolor{dicempleo}{HTML}{1B7F6B}
\definecolor{dichoras}{HTML}{86CBB8}
\definecolor{dicactividad}{HTML}{0F5247}
\definecolor{dicparo}{HTML}{A23B2A}
\definecolor{dicproductividad}{HTML}{B8860B}
% Hogares: renta disponible (pizarra) y tasa de ahorro (negro); el consumo es
% dicconspriv.
\definecolor{dicrenta}{HTML}{4A5A6A}
\definecolor{dicahorro}{HTML}{1A1A1A}
% Política monetaria: tipos del BCE en azules (depósito, el de la zona del
% euro; operaciones principales, medio; facilidad marginal, claro); balance:
% préstamos a los bancos (dorado claro), resto del activo (gris claro).
\definecolor{dicdeposito}{HTML}{17406D}
\definecolor{dicopf}{HTML}{4F86C6}
\definecolor{dicfmc}{HTML}{9EC5EA}
\definecolor{dicprestamos}{HTML}{D9A520}
\definecolor{dicresto}{HTML}{CFCFCF}
% Ciclo: series adelantadas, coincidentes (el PIB) y retrasadas
\definecolor{dicadelantada}{HTML}{D55E00}
\definecolor{dicretrasada}{HTML}{6AA6DA}
% Referencias: rejilla, previsión, sombreados
\definecolor{dicreferencia}{HTML}{767171}

\tikzset{
  % barra de un componente: \addplot[dicbarra=dicfbcf]
  dicbarra/.style={fill=#1, draw=#1!65!black, line width=0.3pt},
  % existencias: gris con trama, para que se distinga también en gris
  dicbarratrama/.style={fill=#1, draw=#1!65!black, line width=0.3pt,
    postaction={pattern=north east lines, pattern color=white}},
  % línea de una serie; con halo blanco cuando pasa sobre barras
  diclinea/.style={very thick, #1},
  diclineahalo/.style={very thick, #1, preaction={draw=white, line width=3.4pt, solid}},
  % puntos de una serie (capa aparte, después de la línea)
  dicpuntos/.style={only marks, mark=*, mark size=1.8pt,
    mark options={fill=white, draw=#1, line width=0.9pt, solid}},
  % línea del cero o de un umbral, después de las barras
  diccero/.style={thin, black},
  dicumbral/.style={thin, dashed, dicreferencia},
}
% Capas de los gráficos de datos (CRITERIOS §7, «tres reglas»), de abajo
% arriba: cuadrícula → sombreados (dicfondo) → barras y áreas (dicbarras) → ejes, cero y umbrales
% (axis lines) → líneas (main) → puntos (dicpuntos) → rótulos (axis
% descriptions). No basta el orden del código: pgfplots dibuja cuadrícula,
% ejes y marcas en sus propias capas (y area style pone los ejes encima), y
% dentro de un axis los \draw, \fill y \node no admiten capa (van con los
% datos de main); por eso todo lo que se dibuja dentro del axis es un \addplot.
% Uso: \begin{axis}[…, dicapas] (al final de las opciones);
%   barras y áreas: \addplot[on layer=dicbarras, …];
%   sombreados:     \dicsombra[estilo]{x1}{y1}{x2}{y2} (rectángulo, en dicfondo,
%                   bajo las barras y sobre la cuadrícula);
%   cero y umbrales: \dicraya[estilo]{(x1,y1) (x2,y2)} (en axis lines);
%   líneas: \addplot normal (main); puntos: \addplot[on layer=dicpuntos, …]
%   (las marcas de cualquier serie van solas a dicpuntos);
%   rótulos: \dicrotulo[opciones del nodo]{(x,y)}{texto} (en axis descriptions),
%   siempre sin fondo (nada de fill=white: si chocan con una línea, se desplazan;
%   CRITERIOS §7, regla 4).
% Los \capa{n} del vídeo no cambian: siguen envolviendo instrucciones completas.
\pgfdeclarelayer{dicfondo}
\pgfdeclarelayer{dicbarras}
\pgfdeclarelayer{dicpuntos}
\pgfplotsset{
  % (el segundo argumento reparte cuadrícula, ejes y rótulos como el estándar)
  layers/dic/.define layer set={axis background, axis grid, dicfondo, dicbarras,
    axis ticks, axis lines, axis tick labels, pre main, main, dicpuntos,
    axis descriptions, axis foreground}{/pgfplots/layers/standard},
  % leyenda sin caja blanca: toma el fondo del gráfico, también si el gráfico va
  % dentro de un cuadro gris (CRITERIOS §7, regla 4, 26/10/2026)
  dicapas/.style={axis on top=false, set layers=dic, mark layer=dicpuntos,
    legend style={fill=none}},
  dicsuelto/.style={forget plot, stack plots=false, sharp plot, mark=none},
}
\newcommand{\dicraya}[2][]{%
  \addplot[on layer=axis lines, dicsuelto, diccero, #1] coordinates {#2};}
\newcommand{\dicsombra}[5][]{%
  \addplot[on layer=dicfondo, dicsuelto, draw=none, #1]
    coordinates {(#2,#3) (#4,#3) (#4,#5) (#2,#5)};}
\newcommand{\dicrotulo}[3][]{%
  \addplot[on layer=axis descriptions, dicsuelto, draw=none] coordinates {#2}
    node[pos=0, #1] {#3};}

%   \begin{document}
%   \begin{tikzpicture}[tcee]
%     \ejes{Q}{P}{6}{5}
%     \draw[curva1] (0.5,4.5) -- (5,0.8) node[etiqueta,right] {$D$};
%     \capa{2}{\draw[curva2] ...;}   % en el vídeo aparece en el paso 2
%   \end{tikzpicture}
%   \end{document}
%
% Paleta: la de la web (morado, dorado, salvia) más tres colores de apoyo
% distinguibles en blanco y negro por el trazo.
% ==========================================================================
\usepackage[T1]{fontenc}
\usepackage{newpxtext,newpxmath}
\usepackage{xcolor}
\usepackage{tikz,pgfplots}
\pgfplotsset{compat=1.18}
\usetikzlibrary{arrows.meta,calc,positioning,patterns,decorations.pathreplacing,intersections,shapes.misc,backgrounds}

\definecolor{tceemorado}{HTML}{5F2987}
\definecolor{tceedorado}{HTML}{B8860B}
\definecolor{tceeverde}{HTML}{3C7D3A}
\definecolor{tceeazul}{HTML}{1F5F99}
\definecolor{tceerojo}{HTML}{B22222}
\definecolor{tceenaranja}{HTML}{E07B28}
\definecolor{tceegris}{HTML}{767171}
\definecolor{tceesalvia}{HTML}{E2EFD9}

% Capas para el vídeo: \capa{n}{…} dibuja su contenido solo a partir del paso n.
% En el PDF, la web y Word se ve todo; generar-video.py compila el gráfico una
% vez por paso con \def\tceepaso{k} para que cada elemento aparezca al narrarlo.
% Vectores con flecha encima, como en el texto (tcee.sty): \vect{x}, \vect{p}
\providecommand{\vect}[1]{\vec{#1}}
\providecommand{\tceepaso}{1000}
\newcommand{\capa}[2]{\ifnum#1>\tceepaso\relax\else#2\fi}
% Gráficos por pasos del mismo tamaño (propuesta del 3.A.29, 11/11/2026): con
% \def\tceefantasma{1} antes del \input (p. ej., x-paso1.tex:
% «\def\tceepaso{3}\def\tceefantasma{1}\input{x.tex}»), las capas posteriores
% al paso se dibujan invisibles y ocupan su sitio. Sin \tceefantasma (el vídeo),
% no cambia nada.
\ifdefined\tceefantasma
  \renewcommand{\capa}[2]{\ifnum#1>\tceepaso\relax\begin{scope}[opacity=0]#2\end{scope}\else#2\fi}
\fi

\tikzset{
  tcee/.style={line cap=round,line join=round,font=\small,>={Stealth[length=2.2mm]}},
  eje/.style={->,thick,black},
  curva1/.style={very thick,tceemorado},
  curva2/.style={very thick,tceeazul},
  curva3/.style={very thick,tceeverde},
  curva4/.style={very thick,tceedorado},
  curvanueva/.style={very thick,tceerojo},
  desplazada/.style={very thick,dashed},
  guia/.style={thin,dashed,tceegris},
  punto/.style={circle,fill=black,inner sep=1.3pt},
  etiqueta/.style={font=\small,inner sep=2pt},
  flecha/.style={->,thick,tceerojo},
  area/.style={fill=tceemorado,fill opacity=0.18,draw=none},
  area2/.style={fill=tceedorado,fill opacity=0.22,draw=none},
}

% \ejes{etiqueta x}{etiqueta y}{largo x}{alto y}
\newcommand{\ejes}[4]{%
  \draw[eje] (0,0) -- (#3,0) node[below left=2pt and 0pt] {#1};
  \draw[eje] (0,0) -- (0,#4) node[left] {#2};
  \node[below left] at (0,0) {$0$};}
% \proyeccion{(x,y)}{etq x}{etq y}: líneas guía a los ejes con etiquetas
\newcommand{\proyeccion}[3]{%
  \draw[guia] let \p1=#1 in (\x1,0) node[below] {#2} |- (0,\y1) node[left] {#3};}

% --------------------------------------------------------------------------
% Paleta «dictamen» (dic…): gráficos de datos de la Guía del dictamen (1D01) y
% de cualquier gráfico de coyuntura que quiera los mismos colores. Un concepto,
% un color en todos los gráficos (CRITERIOS §7, «tres reglas», 24/10/2026);
% tabla completa en fuentes/1D01/_trabajo/paleta-1D01.md. No cambia nada de
% lo anterior: son colores y estilos nuevos.
% Actividad: el PIB (y sus variantes: PIB por persona, PIB mundial), morado.
\definecolor{dicpib}{HTML}{5F2987}
% Demanda interna, gama cálida: consumo (naranjas: privado oscuro, público
% claro), inversión (verde) y existencias (gris con trama, aparte).
\definecolor{dicconspriv}{HTML}{C76A1F}
\definecolor{dicconspub}{HTML}{F5C690}
\definecolor{dicfbcf}{HTML}{2A6B3C}
\definecolor{dicexistencias}{HTML}{B4B4B4}
% Demanda externa y comercio, gama fría (azul claro, que se separa del morado)
\definecolor{dicexterna}{HTML}{6AA6DA}
% Zonas: zona del euro (azul marino), España (rojo), resto de países (grises)
\definecolor{dicze}{HTML}{17406D}
\definecolor{dicespana}{HTML}{C8382E}
\definecolor{dicotro}{HTML}{8A8A8A}
\definecolor{dicotroclaro}{HTML}{B5B5B5}
\definecolor{dicotromasclaro}{HTML}{DADADA}
% Precios: inflación general (negro, continuo) y subyacente (negro,
% discontinuo); componentes: energía (bermellón), alimentos (amarillo) y los
% de la subyacente en azules (bienes industriales claro, servicios oscuro).
\definecolor{dicinflacion}{HTML}{1A1A1A}
\definecolor{dicenergia}{HTML}{D55E00}
\definecolor{dicalimentos}{HTML}{F2D45C}
\definecolor{dicbine}{HTML}{7FB2E0}
\definecolor{dicservicios}{HTML}{1F4E89}
% Trabajo: empleo (verde azulado), horas (el mismo, claro), tasa de actividad
% (verde azulado oscuro), paro (teja); productividad (dorado de la web).
\definecolor{dicempleo}{HTML}{1B7F6B}
\definecolor{dichoras}{HTML}{86CBB8}
\definecolor{dicactividad}{HTML}{0F5247}
\definecolor{dicparo}{HTML}{A23B2A}
\definecolor{dicproductividad}{HTML}{B8860B}
% Hogares: renta disponible (pizarra) y tasa de ahorro (negro); el consumo es
% dicconspriv.
\definecolor{dicrenta}{HTML}{4A5A6A}
\definecolor{dicahorro}{HTML}{1A1A1A}
% Política monetaria: tipos del BCE en azules (depósito, el de la zona del
% euro; operaciones principales, medio; facilidad marginal, claro); balance:
% préstamos a los bancos (dorado claro), resto del activo (gris claro).
\definecolor{dicdeposito}{HTML}{17406D}
\definecolor{dicopf}{HTML}{4F86C6}
\definecolor{dicfmc}{HTML}{9EC5EA}
\definecolor{dicprestamos}{HTML}{D9A520}
\definecolor{dicresto}{HTML}{CFCFCF}
% Ciclo: series adelantadas, coincidentes (el PIB) y retrasadas
\definecolor{dicadelantada}{HTML}{D55E00}
\definecolor{dicretrasada}{HTML}{6AA6DA}
% Referencias: rejilla, previsión, sombreados
\definecolor{dicreferencia}{HTML}{767171}

\tikzset{
  % barra de un componente: \addplot[dicbarra=dicfbcf]
  dicbarra/.style={fill=#1, draw=#1!65!black, line width=0.3pt},
  % existencias: gris con trama, para que se distinga también en gris
  dicbarratrama/.style={fill=#1, draw=#1!65!black, line width=0.3pt,
    postaction={pattern=north east lines, pattern color=white}},
  % línea de una serie; con halo blanco cuando pasa sobre barras
  diclinea/.style={very thick, #1},
  diclineahalo/.style={very thick, #1, preaction={draw=white, line width=3.4pt, solid}},
  % puntos de una serie (capa aparte, después de la línea)
  dicpuntos/.style={only marks, mark=*, mark size=1.8pt,
    mark options={fill=white, draw=#1, line width=0.9pt, solid}},
  % línea del cero o de un umbral, después de las barras
  diccero/.style={thin, black},
  dicumbral/.style={thin, dashed, dicreferencia},
}
% Capas de los gráficos de datos (CRITERIOS §7, «tres reglas»), de abajo
% arriba: cuadrícula → sombreados (dicfondo) → barras y áreas (dicbarras) → ejes, cero y umbrales
% (axis lines) → líneas (main) → puntos (dicpuntos) → rótulos (axis
% descriptions). No basta el orden del código: pgfplots dibuja cuadrícula,
% ejes y marcas en sus propias capas (y area style pone los ejes encima), y
% dentro de un axis los \draw, \fill y \node no admiten capa (van con los
% datos de main); por eso todo lo que se dibuja dentro del axis es un \addplot.
% Uso: \begin{axis}[…, dicapas] (al final de las opciones);
%   barras y áreas: \addplot[on layer=dicbarras, …];
%   sombreados:     \dicsombra[estilo]{x1}{y1}{x2}{y2} (rectángulo, en dicfondo,
%                   bajo las barras y sobre la cuadrícula);
%   cero y umbrales: \dicraya[estilo]{(x1,y1) (x2,y2)} (en axis lines);
%   líneas: \addplot normal (main); puntos: \addplot[on layer=dicpuntos, …]
%   (las marcas de cualquier serie van solas a dicpuntos);
%   rótulos: \dicrotulo[opciones del nodo]{(x,y)}{texto} (en axis descriptions),
%   siempre sin fondo (nada de fill=white: si chocan con una línea, se desplazan;
%   CRITERIOS §7, regla 4).
% Los \capa{n} del vídeo no cambian: siguen envolviendo instrucciones completas.
\pgfdeclarelayer{dicfondo}
\pgfdeclarelayer{dicbarras}
\pgfdeclarelayer{dicpuntos}
\pgfplotsset{
  % (el segundo argumento reparte cuadrícula, ejes y rótulos como el estándar)
  layers/dic/.define layer set={axis background, axis grid, dicfondo, dicbarras,
    axis ticks, axis lines, axis tick labels, pre main, main, dicpuntos,
    axis descriptions, axis foreground}{/pgfplots/layers/standard},
  % leyenda sin caja blanca: toma el fondo del gráfico, también si el gráfico va
  % dentro de un cuadro gris (CRITERIOS §7, regla 4, 26/10/2026)
  dicapas/.style={axis on top=false, set layers=dic, mark layer=dicpuntos,
    legend style={fill=none}},
  dicsuelto/.style={forget plot, stack plots=false, sharp plot, mark=none},
}
\newcommand{\dicraya}[2][]{%
  \addplot[on layer=axis lines, dicsuelto, diccero, #1] coordinates {#2};}
\newcommand{\dicsombra}[5][]{%
  \addplot[on layer=dicfondo, dicsuelto, draw=none, #1]
    coordinates {(#2,#3) (#4,#3) (#4,#5) (#2,#5)};}
\newcommand{\dicrotulo}[3][]{%
  \addplot[on layer=axis descriptions, dicsuelto, draw=none] coordinates {#2}
    node[pos=0, #1] {#3};}

%   \begin{document}
%   \begin{tikzpicture}[tcee]
%     \ejes{Q}{P}{6}{5}
%     \draw[curva1] (0.5,4.5) -- (5,0.8) node[etiqueta,right] {$D$};
%     \capa{2}{\draw[curva2] ...;}   % en el vídeo aparece en el paso 2
%   \end{tikzpicture}
%   \end{document}
%
% Paleta: la de la web (morado, dorado, salvia) más tres colores de apoyo
% distinguibles en blanco y negro por el trazo.
% ==========================================================================
\usepackage[T1]{fontenc}
\usepackage{newpxtext,newpxmath}
\usepackage{xcolor}
\usepackage{tikz,pgfplots}
\pgfplotsset{compat=1.18}
\usetikzlibrary{arrows.meta,calc,positioning,patterns,decorations.pathreplacing,intersections,shapes.misc,backgrounds}

\definecolor{tceemorado}{HTML}{5F2987}
\definecolor{tceedorado}{HTML}{B8860B}
\definecolor{tceeverde}{HTML}{3C7D3A}
\definecolor{tceeazul}{HTML}{1F5F99}
\definecolor{tceerojo}{HTML}{B22222}
\definecolor{tceenaranja}{HTML}{E07B28}
\definecolor{tceegris}{HTML}{767171}
\definecolor{tceesalvia}{HTML}{E2EFD9}

% Capas para el vídeo: \capa{n}{…} dibuja su contenido solo a partir del paso n.
% En el PDF, la web y Word se ve todo; generar-video.py compila el gráfico una
% vez por paso con \def\tceepaso{k} para que cada elemento aparezca al narrarlo.
% Vectores con flecha encima, como en el texto (tcee.sty): \vect{x}, \vect{p}
\providecommand{\vect}[1]{\vec{#1}}
\providecommand{\tceepaso}{1000}
\newcommand{\capa}[2]{\ifnum#1>\tceepaso\relax\else#2\fi}
% Gráficos por pasos del mismo tamaño (propuesta del 3.A.29, 11/11/2026): con
% \def\tceefantasma{1} antes del \input (p. ej., x-paso1.tex:
% «\def\tceepaso{3}\def\tceefantasma{1}\input{x.tex}»), las capas posteriores
% al paso se dibujan invisibles y ocupan su sitio. Sin \tceefantasma (el vídeo),
% no cambia nada.
\ifdefined\tceefantasma
  \renewcommand{\capa}[2]{\ifnum#1>\tceepaso\relax\begin{scope}[opacity=0]#2\end{scope}\else#2\fi}
\fi

\tikzset{
  tcee/.style={line cap=round,line join=round,font=\small,>={Stealth[length=2.2mm]}},
  eje/.style={->,thick,black},
  curva1/.style={very thick,tceemorado},
  curva2/.style={very thick,tceeazul},
  curva3/.style={very thick,tceeverde},
  curva4/.style={very thick,tceedorado},
  curvanueva/.style={very thick,tceerojo},
  desplazada/.style={very thick,dashed},
  guia/.style={thin,dashed,tceegris},
  punto/.style={circle,fill=black,inner sep=1.3pt},
  etiqueta/.style={font=\small,inner sep=2pt},
  flecha/.style={->,thick,tceerojo},
  area/.style={fill=tceemorado,fill opacity=0.18,draw=none},
  area2/.style={fill=tceedorado,fill opacity=0.22,draw=none},
}

% \ejes{etiqueta x}{etiqueta y}{largo x}{alto y}
\newcommand{\ejes}[4]{%
  \draw[eje] (0,0) -- (#3,0) node[below left=2pt and 0pt] {#1};
  \draw[eje] (0,0) -- (0,#4) node[left] {#2};
  \node[below left] at (0,0) {$0$};}
% \proyeccion{(x,y)}{etq x}{etq y}: líneas guía a los ejes con etiquetas
\newcommand{\proyeccion}[3]{%
  \draw[guia] let \p1=#1 in (\x1,0) node[below] {#2} |- (0,\y1) node[left] {#3};}

% --------------------------------------------------------------------------
% Paleta «dictamen» (dic…): gráficos de datos de la Guía del dictamen (1D01) y
% de cualquier gráfico de coyuntura que quiera los mismos colores. Un concepto,
% un color en todos los gráficos (CRITERIOS §7, «tres reglas», 24/10/2026);
% tabla completa en fuentes/1D01/_trabajo/paleta-1D01.md. No cambia nada de
% lo anterior: son colores y estilos nuevos.
% Actividad: el PIB (y sus variantes: PIB por persona, PIB mundial), morado.
\definecolor{dicpib}{HTML}{5F2987}
% Demanda interna, gama cálida: consumo (naranjas: privado oscuro, público
% claro), inversión (verde) y existencias (gris con trama, aparte).
\definecolor{dicconspriv}{HTML}{C76A1F}
\definecolor{dicconspub}{HTML}{F5C690}
\definecolor{dicfbcf}{HTML}{2A6B3C}
\definecolor{dicexistencias}{HTML}{B4B4B4}
% Demanda externa y comercio, gama fría (azul claro, que se separa del morado)
\definecolor{dicexterna}{HTML}{6AA6DA}
% Zonas: zona del euro (azul marino), España (rojo), resto de países (grises)
\definecolor{dicze}{HTML}{17406D}
\definecolor{dicespana}{HTML}{C8382E}
\definecolor{dicotro}{HTML}{8A8A8A}
\definecolor{dicotroclaro}{HTML}{B5B5B5}
\definecolor{dicotromasclaro}{HTML}{DADADA}
% Precios: inflación general (negro, continuo) y subyacente (negro,
% discontinuo); componentes: energía (bermellón), alimentos (amarillo) y los
% de la subyacente en azules (bienes industriales claro, servicios oscuro).
\definecolor{dicinflacion}{HTML}{1A1A1A}
\definecolor{dicenergia}{HTML}{D55E00}
\definecolor{dicalimentos}{HTML}{F2D45C}
\definecolor{dicbine}{HTML}{7FB2E0}
\definecolor{dicservicios}{HTML}{1F4E89}
% Trabajo: empleo (verde azulado), horas (el mismo, claro), tasa de actividad
% (verde azulado oscuro), paro (teja); productividad (dorado de la web).
\definecolor{dicempleo}{HTML}{1B7F6B}
\definecolor{dichoras}{HTML}{86CBB8}
\definecolor{dicactividad}{HTML}{0F5247}
\definecolor{dicparo}{HTML}{A23B2A}
\definecolor{dicproductividad}{HTML}{B8860B}
% Hogares: renta disponible (pizarra) y tasa de ahorro (negro); el consumo es
% dicconspriv.
\definecolor{dicrenta}{HTML}{4A5A6A}
\definecolor{dicahorro}{HTML}{1A1A1A}
% Política monetaria: tipos del BCE en azules (depósito, el de la zona del
% euro; operaciones principales, medio; facilidad marginal, claro); balance:
% préstamos a los bancos (dorado claro), resto del activo (gris claro).
\definecolor{dicdeposito}{HTML}{17406D}
\definecolor{dicopf}{HTML}{4F86C6}
\definecolor{dicfmc}{HTML}{9EC5EA}
\definecolor{dicprestamos}{HTML}{D9A520}
\definecolor{dicresto}{HTML}{CFCFCF}
% Ciclo: series adelantadas, coincidentes (el PIB) y retrasadas
\definecolor{dicadelantada}{HTML}{D55E00}
\definecolor{dicretrasada}{HTML}{6AA6DA}
% Referencias: rejilla, previsión, sombreados
\definecolor{dicreferencia}{HTML}{767171}

\tikzset{
  % barra de un componente: \addplot[dicbarra=dicfbcf]
  dicbarra/.style={fill=#1, draw=#1!65!black, line width=0.3pt},
  % existencias: gris con trama, para que se distinga también en gris
  dicbarratrama/.style={fill=#1, draw=#1!65!black, line width=0.3pt,
    postaction={pattern=north east lines, pattern color=white}},
  % línea de una serie; con halo blanco cuando pasa sobre barras
  diclinea/.style={very thick, #1},
  diclineahalo/.style={very thick, #1, preaction={draw=white, line width=3.4pt, solid}},
  % puntos de una serie (capa aparte, después de la línea)
  dicpuntos/.style={only marks, mark=*, mark size=1.8pt,
    mark options={fill=white, draw=#1, line width=0.9pt, solid}},
  % línea del cero o de un umbral, después de las barras
  diccero/.style={thin, black},
  dicumbral/.style={thin, dashed, dicreferencia},
}
% Capas de los gráficos de datos (CRITERIOS §7, «tres reglas»), de abajo
% arriba: cuadrícula → sombreados (dicfondo) → barras y áreas (dicbarras) → ejes, cero y umbrales
% (axis lines) → líneas (main) → puntos (dicpuntos) → rótulos (axis
% descriptions). No basta el orden del código: pgfplots dibuja cuadrícula,
% ejes y marcas en sus propias capas (y area style pone los ejes encima), y
% dentro de un axis los \draw, \fill y \node no admiten capa (van con los
% datos de main); por eso todo lo que se dibuja dentro del axis es un \addplot.
% Uso: \begin{axis}[…, dicapas] (al final de las opciones);
%   barras y áreas: \addplot[on layer=dicbarras, …];
%   sombreados:     \dicsombra[estilo]{x1}{y1}{x2}{y2} (rectángulo, en dicfondo,
%                   bajo las barras y sobre la cuadrícula);
%   cero y umbrales: \dicraya[estilo]{(x1,y1) (x2,y2)} (en axis lines);
%   líneas: \addplot normal (main); puntos: \addplot[on layer=dicpuntos, …]
%   (las marcas de cualquier serie van solas a dicpuntos);
%   rótulos: \dicrotulo[opciones del nodo]{(x,y)}{texto} (en axis descriptions),
%   siempre sin fondo (nada de fill=white: si chocan con una línea, se desplazan;
%   CRITERIOS §7, regla 4).
% Los \capa{n} del vídeo no cambian: siguen envolviendo instrucciones completas.
\pgfdeclarelayer{dicfondo}
\pgfdeclarelayer{dicbarras}
\pgfdeclarelayer{dicpuntos}
\pgfplotsset{
  % (el segundo argumento reparte cuadrícula, ejes y rótulos como el estándar)
  layers/dic/.define layer set={axis background, axis grid, dicfondo, dicbarras,
    axis ticks, axis lines, axis tick labels, pre main, main, dicpuntos,
    axis descriptions, axis foreground}{/pgfplots/layers/standard},
  % leyenda sin caja blanca: toma el fondo del gráfico, también si el gráfico va
  % dentro de un cuadro gris (CRITERIOS §7, regla 4, 26/10/2026)
  dicapas/.style={axis on top=false, set layers=dic, mark layer=dicpuntos,
    legend style={fill=none}},
  dicsuelto/.style={forget plot, stack plots=false, sharp plot, mark=none},
}
\newcommand{\dicraya}[2][]{%
  \addplot[on layer=axis lines, dicsuelto, diccero, #1] coordinates {#2};}
\newcommand{\dicsombra}[5][]{%
  \addplot[on layer=dicfondo, dicsuelto, draw=none, #1]
    coordinates {(#2,#3) (#4,#3) (#4,#5) (#2,#5)};}
\newcommand{\dicrotulo}[3][]{%
  \addplot[on layer=axis descriptions, dicsuelto, draw=none] coordinates {#2}
    node[pos=0, #1] {#3};}

\begin{document}
\hyphenpenalty=10000 \exhyphenpenalty=10000
% Compraventa (merchanting) en la balanza de pagos (MBP6, párr. 10.41-10.49):
% el comerciante residente en A compra la mercancía a B por 80 y la vende a C
% por 100; la mercancía va de B a C sin pasar por A. En la balanza de A se
% anota la exportación neta de bienes en compraventa: +20.
% Trazos: mercancía en azul continuo (curva2), propiedad en morado discontinuo
% (curva1), pagos en dorado punteado (curva4).
% Capas: los 3 países -> la mercancía -> la compra de A (propiedad y pago) ->
% la venta de A (propiedad y pago) -> el asiento en la balanza de A.
\begin{tikzpicture}[tcee, font=\footnotesize,
  pais/.style={draw=#1, fill=#1!12, line width=0.8pt, rounded corners=3pt,
    align=center, text width=2.9cm, minimum height=1.25cm, inner sep=3pt},
  merc/.style={curva2, -{Stealth[length=2.6mm]}},
  prop/.style={curva1, dashed, -{Stealth[length=2.6mm]}},
  pago/.style={curva4, densely dotted, -{Stealth[length=2.6mm]}},
  rot/.style={font=\scriptsize, align=center, inner sep=2pt}]
  % Marco fijo: todas las capas del vídeo con el mismo tamaño
  \path (-6.1,-2.3) rectangle (6.1,3.55);
  % 1: los 3 países
  \capa{1}{\node[pais=tceegris] (B) at (-4.3,0) {\textbf{País B}\\ productor};
    \node[pais=tceegris] (C) at (4.3,0) {\textbf{País C}\\ comprador};
    \node[pais=tceemorado, fill=tceesalvia, text width=3.5cm] (A) at (0,2.7) {\textbf{País A}\\ comerciante residente\\ (economía que compila)};}
  % 2: la mercancía va de B a C sin pasar por A
  \capa{2}{\draw[merc] (B.east) -- (C.west)
    node[midway, below, rot, text=tceeazul] {la mercancía va de B a C\\ sin pasar por A};}
  % 3: A compra a B: le pasa la propiedad...
  \capa{3}{\draw[prop] ([xshift=6mm]B.north) to[bend left=12]
      node[pos=0.45, above left=-1pt, rot, text=tceemorado] {propiedad:\\ A compra por 80} ([yshift=-1mm]A.west);}
  % 4: ... y A le paga 80
  \capa{4}{\draw[pago] ([yshift=-5mm]A.west) to[bend left=12]
      node[pos=0.5, below right=-1pt, rot, text=tceedorado!80!black] {pago: 80} ([xshift=13mm]B.north);}
  % 5: A vende a C: le pasa la propiedad...
  \capa{5}{\draw[prop] ([yshift=-1mm]A.east) to[bend left=12]
      node[pos=0.55, above right=-1pt, rot, text=tceemorado] {propiedad:\\ A vende por 100} ([xshift=-6mm]C.north);}
  % 6: ... y C le paga 100
  \capa{6}{\draw[pago] ([xshift=-13mm]C.north) to[bend left=12]
      node[pos=0.5, below left=-1pt, rot, text=tceedorado!80!black] {pago: 100} ([yshift=-5mm]A.east);}
  % 7: el asiento en la balanza de A
  \capa{7}{\node[draw=tceemorado, rounded corners=2pt, align=center, inner sep=4pt] at (0,-1.75)
    {En la balanza de A: \textbf{exportación neta de bienes en compraventa} $=100-80=+20$};}
\end{tikzpicture}
\end{document}
